
\chapter{2016年北京卷（理）}

\section{单选题}

\begin{problem}
已知集合 \(A=\{x\mid |x|<2\}\)，\(B=\{-1,0,1,2,3\}\)，则 \(A\cap B=\)（\quad）

\choices
    {\(\{0,1\}\)}
    {\(\{0,1,2\}\)}
    {\(\{-1,0,1\}\)}
    {\(\{-1,0,1,2\}\)}
\end{problem}

\begin{problem}
若 \(x\)，\(y\) 满足 \(\begin{cases}
2x-y\le 0,\\
x+y\le 3,\\
x\ge 0,
\end{cases}\)
则 \(2x+y\) 的最大值为（\quad）

\choices
    {\(0\)}
    {\(3\)}
    {\(4\)}
    {\(5\)}
\end{problem}

\begin{problem}
执行如图所示的程序框图，若输入的 \(a\) 值为 \(1\)，则输出的 \(k\) 值为（\quad）

\begingroup
\bitmapfigure[max width=6.0cm]{img/2016/beijing_science/q3_fig1.png}
\endgroup

\choices
    {\(1\)}
    {\(2\)}
    {\(3\)}
    {\(4\)}
\end{problem}

\begin{problem}
设 \(\bs a\)，\(\bs b\) 是向量，则“\(|\bs a|=|\bs b|\)”是“\(|\bs a+\bs b|=|\bs a-\bs b|\)”的（\quad）

\choices
    {充分而不必要条件}
    {必要而不充分条件}
    {充分必要条件}
    {既不充分也不必要条件}
\end{problem}

\begin{problem}
已知 \(x,y\in\R\)，且 \(x>y>0\)，则（\quad）

\choices
    {\(\dfrac1x-\dfrac1y>0\)}
    {\(\sin x-\sin y>0\)}
    {\(\left(\dfrac12\right)^x-\left(\dfrac12\right)^y<0\)}
    {\(\ln x+\ln y>0\)}
\end{problem}

\begin{problem}
某三棱锥的三视图如图所示，则该三棱锥的体积为（\quad）

\begingroup
\bitmapfigure[max width=0.48\linewidth]{img/2016/beijing_science/q6_fig1.png}
\endgroup

\choices
    {\(\dfrac16\)}
    {\(\dfrac13\)}
    {\(\dfrac12\)}
    {\(1\)}
\end{problem}

\begin{problem}
将函数 \(y=\sin\left(2x-\dfrac{\pi}{3}\right)\) 图象上的点 \(P\left(\dfrac{\pi}{4},t\right)\) 向左平移 \(s\)（\(s>0\)）个单位长度得到点 \(P'\)，若 \(P'\) 位于函数 \(y=\sin2x\) 的图象上，则（\quad）

\choices
    {\(t=\dfrac12\)，\(s\) 的最小值为 \(\dfrac{\pi}{6}\)}
    {\(t=\dfrac{\sqrt3}{2}\)，\(s\) 的最小值为 \(\dfrac{\pi}{6}\)}
    {\(t=\dfrac12\)，\(s\) 的最小值为 \(\dfrac{\pi}{3}\)}
    {\(t=\dfrac{\sqrt3}{2}\)，\(s\) 的最小值为 \(\dfrac{\pi}{3}\)}
\end{problem}

\begin{problem}
袋中装有偶数个球，其中红球，黑球各占一半. 甲，乙，丙是三个空盒. 每次从袋中任意取出两个球，将其中一个球放入甲盒，如果这个球是红球，就将另一个放入乙盒，否则就放入丙盒. 重复上述过程，直到袋中所有球都被放入盒中，则（\quad）

\choices
    {乙盒中黑球不多于丙盒中黑球}
    {乙盒中红球与丙盒中黑球一样多}
    {乙盒中红球不多于丙盒中红球}
    {乙盒中黑球与丙盒中红球一样多}
\end{problem}

\section{填空题}

\begin{problem}
设 \(a\in\R\)，若复数 \((1+\i)(a+\i)\) 在复平面内对应的点位于实轴上，则 \(a=\fillinblank{}\).
\end{problem}

\begin{problem}
在 \((1-2x)^6\) 的展开式中，\(x^2\) 的系数为\fillinblank{}（用数字作答）.
\end{problem}

\begin{problem}
在极坐标系中，直线 \(\rho\cos\theta-\sqrt{3}\rho\sin\theta-1=0\) 与圆 \(\rho=2\cos\theta\) 交于 \(A\)，\(B\) 两点，则 \(|AB|=\fillinblank{}\).
\end{problem}

\begin{problem}
已知 \(\{a_n\}\) 为等差数列，\(S_n\) 为其前 \(n\) 项和. 若 \(a_1=6\)，\(a_3+a_5=0\)，则 \(S_6=\fillinblank{}\).
\end{problem}

\begin{problem}
双曲线 \(\dfrac{x^2}{a^2}-\dfrac{y^2}{b^2}=1\)（\(a>0\)，\(b>0\)）的渐近线为正方形 \(OABC\) 的边 \(OA\)，\(OC\) 所在的直线，点 \(B\) 为该双曲线的焦点. 若正方形 \(OABC\) 的边长为 \(2\)，则 \(a=\fillinblank{}\).
\end{problem}

\begin{problem}
设函数
\(f(x)=\begin{cases}
x^3-3x, & x\le a,\\
-2x, & x>a.
\end{cases}\)

\begin{enumerate}
\item 若 \(a=0\)，则 \(f(x)\) 的最大值为\fillinblank{}；
\item 若 \(f(x)\) 无最大值，则实数 \(a\) 的取值范围是\fillinblank{}.
\end{enumerate}
\end{problem}

\section{解答题}

\begin{problem}
在 \(\triangle ABC\) 中，\(a^2+c^2=b^2+\sqrt{2}ac\).

\begin{enumerate}
\item 求 \(\angle B\) 的大小；
\item 求 \(\sqrt{2}\cos A+\cos C\) 的最大值.
\end{enumerate}
\end{problem}

\begin{problem}
\(A\)，\(B\)，\(C\) 三个班共有 \(100\) 名学生，为调查他们的体育锻炼情况，通过分层抽样获得了部分学生一周的锻炼时间，数据如表（单位：小时）：

\begin{center}
\begin{tabular}{|c|cccccccc|}
\hline
\(A\)班 & \(6\) & \(6.5\) & \(7\) & \(7.5\) & \(8\) & & & \\
\hline
\(B\)班 & \(6\) & \(7\) & \(8\) & \(9\) & \(10\) & \(11\) & \(12\) & \\
\hline
\(C\)班 & \(3\) & \(4.5\) & \(6\) & \(7.5\) & \(9\) & \(10.5\) & \(12\) & \(13.5\) \\
\hline
\end{tabular}
\end{center}

\begin{enumerate}
\item 试估计 \(C\) 班的学生人数；
\item 从 A 班和 C 班抽出的学生中，各随机选取一个人，A 班选出的人记为甲，C 班选出的人记为乙. 假设所有学生的锻炼时间相对独立，求该周甲的锻炼时间比乙的锻炼时间长的概率；
\item 再从 A，B，C 三班中各随机抽取一名学生，他们该周锻炼时间分别是 \(7\)，\(9\)，\(8.25\)（单位：小时），这 \(3\) 个新数据与表格中的数据构成的新样本的平均数记为 \(\mu_1\)，表格中数据的平均数记为 \(\mu_0\)，试判断 \(\mu_0\) 和 \(\mu_1\) 的大小（结论不要求证明）.
\end{enumerate}
\end{problem}

\begin{problem}
如图，在四棱锥 \(P-ABCD\) 中，平面 \(PAD\perp\) 平面 \(ABCD\)，\(PA\perp PD\)，\(PA=PD\)，\(AB\perp AD\)，\(AB=1\)，\(AD=2\)，\(AC=CD=\sqrt5\).

\begin{enumerate}
\item 求证：\(PD\perp\) 平面 \(PAB\)；
\item 求直线 \(PB\) 与平面 \(PCD\) 所成角的正弦值；
\item 在棱 \(PA\) 上是否存在点 \(M\)，使得 \(BM\parallel\) 平面 \(PCD\)？若存在，求 \(\dfrac{AM}{AP}\) 的值，若不存在，说明理由.
\end{enumerate}

\begingroup
\bitmapfigure[float=inline,max width=0.55\linewidth]{img/2016/beijing_science/q17_fig2.png}
\endgroup
\end{problem}

\begin{problem}
设函数 \(f(x)=x\e^{a-x}+bx\)，曲线 \(y=f(x)\) 在点 \((2,f(2))\) 处的切线方程为 \(y=(\e-1)x+4\).

\begin{enumerate}
\item 求 \(a\)，\(b\) 的值；
\item 求 \(f(x)\) 的单调区间.
\end{enumerate}
\end{problem}

\begin{problem}
已知椭圆 \(C:\dfrac{x^2}{a^2}+\dfrac{y^2}{b^2}=1\)（\(a>b>0\)）的离心率为 \(\dfrac{\sqrt3}{2}\)，\(A(a,0)\)，\(B(0,b)\)，\(O(0,0)\)，\(\triangle OAB\) 的面积为 \(1\).

\begin{enumerate}
\item 求椭圆 \(C\) 的方程；
\item 设 \(P\) 是椭圆 \(C\) 上一点，直线 \(PA\) 与 \(y\) 轴交于点 \(M\)，直线 \(PB\) 与 \(x\) 轴交于点 \(N\). 求证：\(|AN|\cdot|BM|\) 为定值.
\end{enumerate}
\end{problem}

\begin{problem}
设数列 \(A\)：\(a_1,a_2,\ldots,a_N\)（\(N\ge2\)）. 如果对小于 \(n\)（\(2\le n\le N\)）的每个正整数 \(k\) 都有 \(a_k<a_n\)，则称 \(n\) 是数列 \(A\) 的一个“\(G\)时刻”，记 \(G(A)\) 是数列 \(A\) 的所有“\(G\)时刻”组成的集合.

\begin{enumerate}
\item 对数列 \(A\)：\(-2\)，\(2\)，\(-1\)，\(1\)，\(3\)，写出 \(G(A)\) 的所有元素；
\item 证明：若数列 \(A\) 中存在 \(a_n\) 使得 \(a_n>a_1\)，则 \(G(A)\ne\varnothing\)；
\item 证明：若数列 \(A\) 满足 \(a_n-a_{n-1}\le1\)（\(n=2,3,\ldots,N\)），则 \(G(A)\) 的元素个数不小于 \(a_N-a_1\).
\end{enumerate}
\end{problem}
